\chapter{積分の計算}\label{章:基本積分}

前章では,定積分を計算するために原始関数を探せばよいことが分かった.そこで今度は,目の前の関数に対して,その原始関数をどう見つけるかを考える.たとえば $x^2$ なら,微分すると $3x^2$ になる $x^3$ を思い出し,係数を直せばよい.ところが $(x+1)^2$ や $\frac{x^2}{x+1}$ は,見た形が違うだけで,同じ方法へ進めることがある.

この章では,まず微分公式を逆向きに読む.そのうえで,展開,割り算,平方完成,部分分数,三角恒等式を使い,知っている原始関数が見える形へ式を整える.新しい積分法を次々に覚える前に,式を変えるだけで何ができるかを確かめよう.

\section{基本公式}\label{節:基本公式の逆}

原始関数を探す出発点は,微分の公式を逆向きに読むことである.本節では,基本的な関数の原始関数をそうして求め,表にまとめる.以後の節では,与えられた関数をこの表のどれかに帰着させることを目標とする.

$F'(x)=f(x)$ ならば,\,$F$ は $f$ の原始関数である.この関係を使うとき,最初から正しい係数まで思い付く必要はない.形の近い候補を微分し,元の関数との差を直せばよい.

$x^r$ で試してみよう.指数を1つ増やした $x^{r+1}$ を微分すると
\[
 \frac{\dd}{\dd x}x^{r+1}=(r+1)x^r
\]
となり,余分な係数 $r+1$ が出る.そこで $r\ne-1$ ならば,両辺を $r+1$ で割って
\[
 \frac{\dd}{\dd x}\frac{x^{r+1}}{r+1}=x^r
\]
を得る.すなわち
\[
 \int x^r\dd x=\frac{x^{r+1}}{r+1}+C
\]
である.指数を1つ増やし,増やした指数で割るのは,微分で指数が1つ減り,その指数が係数として出ることを元に戻しているのである.

ここで $r=-1$ だけは別に考える必要がある.上の式へ $r=-1$ を代入すると分母が $0$ になり,意味を持たない.また,\,$x^0=1$ を微分しても $x^{-1}$ は現れない.代わりに,対数の微分
\[
 \frac{\dd}{\dd x}\log x=\frac1x\qquad(x>0)
\]
を使う.\,$x<0$ では $\log(-x)$ を微分すれば
\[
 \frac{\dd}{\dd x}\log(-x)
 =\frac1{-x}\,(-1)=\frac1x
\]
となるから,2つの形は $\log|x|$ とまとめられる.

ほかの関数も,微分の公式の左右を入れ換えて読めばよい.たとえば $a^x$ を微分すると $\log a$ が掛かるので,積分では $\log a$ で割る.\footnote{$a=1$ を除くのは,\,$\log1=0$ で割ることになるためである.\,$1^x=1$ の積分そのものは $x+C$ である.}以上をまとめると,次の表を得る.表の右の式を微分すれば,左の式に戻ることを確かめられる.

\begin{theorem}{基本的な不定積分}
次の公式は,右辺が微分でき,左辺の関数が定義される区間で成り立つ.
\[
\begin{array}{c|c}
 f(x)&\int f(x)\dd x\\ \hline
 x^r\quad(r\ne-1)&\frac{x^{r+1}}{r+1}+C\\[5pt]
 \frac1x&\log|x|+C\\[5pt]
 \e^x&\e^x+C\\[3pt]
 a^x\quad(a>0,\ a\ne1)&\frac{a^x}{\log a}+C\\[5pt]
 \sin x&-\cos x+C\\[3pt]
 \cos x&\sin x+C\\[3pt]
 \frac1{\cos^2x}&\tan x+C
\end{array}
\]
実数 $r$ に対するべきの公式は,まず $x>0$ で考える.整数のべきなど,負の $x$ でも定義される場合は,その定義域内の区間にも適用できる.
\end{theorem}

また,微分は和と定数倍を保つ.したがって,和は項ごとに積分でき,定数倍は積分の外に出せる.

ここで1つ断りをいれておこう.不定積分の $C$ は,考えている1つの区間の上での任意定数である.\,$\frac1x$ の定義域は $0$ の左右に分かれ,\,$0$ を越えて1つの区間にすることはできない.そのため,左側で加える定数と右側で加える定数が同じである必要はない.図\ref{図:対数の区間}に,この区間の分かれ方を数直線で示した.本章では以後,不定積分の等式は被積分関数が定義される各区間の上で成り立つものとし,任意定数は区間ごとに独立に選べるものとする.どの点を除き,どの区間で考えるかは,例題の解法で述べる.

\begin{figure}[htbp]
 \centering
 \begin{tikzpicture}[x=1cm,y=1cm]
  \draw[-{Stealth}] (-4,0)--(4,0) node[right] {$x$};
  \draw[very thick,blue] (-3.7,0)--(-0.12,0);
  \draw[very thick,blue] (0.12,0)--(3.7,0);
  \filldraw[fill=white] (0,0) circle (2pt);
  \node[below] at (0,-0.08) {$0$};
  \node[above] at (-2,0.18) {$\log(-x)+C_-$};
  \node[above] at (2,0.18) {$\log x+C_+$};
  \node[below] at (-2,-0.35) {$x<0$};
  \node[below] at (2,-0.35) {$x>0$};
 \end{tikzpicture}
 \caption{$0$ を越えて1つの区間にはできない.定数は左右で独立に選べる.}
 \label{図:対数の区間}
\end{figure}

\begin{bookpoint}
\textbf{POINT}\par
根号や分数を $x^r$ の形に書き直し,和と定数倍に分けて,項ごとに表の公式を使う.
\end{bookpoint}

\begin{example}{べきの積分}
$x>0$ において,\,$\int\left(2\sqrt{x}+\frac3{x^2}\right)\dd x$ を求めよ.
\end{example}

\paragraphbf{方針}
$\sqrt{x}=x^{\frac12}$,\,$\frac1{x^2}=x^{-2}$ と書き直し,べきの公式に帰着させる.

\paragraphbf{解答}
$\frac43x^{\frac32}-\frac3x+C$

\paragraphbf{解法}
$\sqrt{x}=x^{\frac12}$,\,$\frac1{x^2}=x^{-2}$ と書き,項ごとに積分すると
\begin{align*}
 \int\left(2\sqrt{x}+\frac3{x^2}\right)\dd x
 &=2\int x^{\frac12}\dd x+3\int x^{-2}\dd x\\
 &=2\cdot\frac23x^{\frac32}+3\cdot\frac{x^{-1}}{-1}+C\\
 &=\frac43x^{\frac32}-\frac3x+C
\end{align*}
となる.\,$x>0$ では両方の項が定義される.分母に $x$ があっても,\,$\frac1{x^2}$ は指数が $-2$ なので,対数ではなくべきの公式を使う.答えを微分すると
\[
 \frac43\cdot\frac32x^{\frac12}+3x^{-2}=2\sqrt{x}+\frac3{x^2}
\]
となり,係数と符号の両方が元の式と一致する.

\paragraphbf{積の積分}
和は項ごとに積分できたが,積について
\[
 \int f(x)g(x)\dd x
 =\left(\int f(x)\dd x\right)
  \left(\int g(x)\dd x\right)
\]
という規則はない.右辺を微分すると積の微分の2項が現れ,ふつうは $f(x)g(x)$ にはならない.実際,\,$f(x)=g(x)=x$ とすると,左辺は $\frac{x^3}3+C$ だが,右辺の $\frac{x^2}2\cdot\frac{x^2}2=\frac{x^4}4$ を微分すると $x^3$ になり,\,$x^2$ には戻らない.積があるときは,まず展開できるか,恒等式で和にできるかを考える.これが次節以降で扱う変形であり,後の章で学ぶ置換積分と部分積分も,この問題への別の答えである.
%
\section{一次式の積分}\label{節:一次式の積分}

前節の表には $\sin x$ や $\e^x$ はあるが,\,$\sin(3x)$ や $\e^{2x-1}$ はない.では,変数 $x$ のところが一次式に置き換わった
\[
 \int f(ax+b)\dd x\qquad(a\ne0)
\]
はどうなるだろうか.\,$f$ の原始関数 $F$ は表から分かっているとして,その $F$ に帰着させることを目標とする.

まずは $\sin(3x)$ で試してみよう.\,$-\cos(3x)$ を候補にして微分すると
\[
 \frac{\dd}{\dd x}\{-\cos(3x)\}=3\sin(3x)
\]
となる.合成関数の微分によって,内側の $3x$ の微分 $3$ が余分に掛かったのだ.したがって候補を $3$ で割ればよい.

一般の場合も同じである.\,$F'=f$ とすると,合成関数の微分から
\[
 \frac{\dd}{\dd x}F(ax+b)=af(ax+b)
\]
となり,余分な $a$ が掛かる.両辺を $a$ で割って
\[
 \frac{\dd}{\dd x}\left\{\frac1aF(ax+b)\right\}=f(ax+b)
\]
を得る.すなわち
\[
 \int f(ax+b)\dd x=\frac1aF(ax+b)+C
\]
である.\,$a$ で割るのは書式上の約束ではなく,微分したときに現れる $a$ を打ち消すためである.

分母が一次式の $\frac1{ax+b}$ では,\,$\log|x|$ から $\frac1a\log|ax+b|+C$ を得る.この場合は $ax+b=0$ となる $x=-\frac ba$ を除き,その左右の各区間で考える.

\begin{bookpoint}
\textbf{POINT}\par
表の原始関数 $F$ に $ax+b$ を代入し,内側の係数 $a$ で割って $\frac1aF(ax+b)$ とする.
\end{bookpoint}

\begin{example}{内側の微分を確かめる}
次の不定積分を求めよ.
\[
 \int\left(4\e^{2x-1}+\frac3{5x+2}-2\cos(3x)\right)\dd x.
\]
\end{example}

\paragraphbf{方針}
3つの項を $\e^x$,\,$\frac1x$,\,$\cos x$ の $x$ を一次式に置き換えたものとみて,内側の係数 $2,5,3$ でそれぞれ割る.

\paragraphbf{解答}
$2\e^{2x-1}+\frac35\log|5x+2|-\frac23\sin(3x)+C$

\paragraphbf{解法}
指数の $2x-1$,分母の $5x+2$,角の $3x$ の $x$ の係数は,それぞれ $2,5,3$ である.これらで割ると
\begin{align*}
 &\int\left(4\e^{2x-1}+\frac3{5x+2}-2\cos(3x)\right)\dd x\\
 &\qquad=4\cdot\frac12\e^{2x-1}+3\cdot\frac15\log|5x+2|-2\cdot\frac13\sin(3x)+C\\
 &\qquad=2\e^{2x-1}+\frac35\log|5x+2|-\frac23\sin(3x)+C
\end{align*}
となる.ただし,\,$5x+2=0$ となる $x=-\frac25$ は除き,\,$x<-\frac25$ と $x>-\frac25$ の各区間で考える.対数の項を微分すると
\[
 \frac35\cdot\frac5{5x+2}=\frac3{5x+2}
\]
となり,分母の係数 $5$ が打ち消されて元の項に戻る.

\paragraphbf{内側が一次式でない場合}
一次式の場合は内側の微分が定数 $a$ なので,係数を直すだけで済んだ.内側が $x^2+1$ ならば,微分は $2x$ であり,定数ではない.\,$F(x^2+1)$ を $2x$ で割った候補を微分しても,商の微分から余分な項が出てしまう.次章の置換積分では,この $2x$ が積分される式の中にあるかどうかが手がかりになる.
%
\section{展開と割り算}\label{節:展開と割算}

$x(2x-1)^2$ や $\frac{x^3+2x^2+3}{x+1}$ は,そのままでは前節までの表のどれにも当てはまらない.しかし,式を同じ値の別の式に書き換えれば,\,$x$ のべきや一次式の逆数の和にできる.本節では,展開と割り算という2つの書き換えによって,基本公式に帰着させる.

ここで1つ断りをいれておこう.割り算や約分で式を書き換えても,積分するのはもとの関数である.たとえば $\frac{x^2-1}{x-1}=x+1$ は $x\ne1$ でしか成り立たない.右辺が $x=1$ でも定義されるからといって,左辺の関数がその点で定義されたことにはならない.したがって,書き換えた後に得た答えも,もとの関数の分母が $0$ にならない各区間で読む.\footnote{積分する区間にその点が含まれる場合は,連続に延長した関数を扱うのか,もとの式を扱うのかを明らかにする必要がある.}

まず整式の積である.積の形では微分の逆が見えなくても,展開すれば各項が $x$ のべきになり,基本公式を1項ずつ使える.式を短く保つことと,積分しやすくすることは同じではないのだ.

次に,整式 $P(x),Q(x)$ の商 $\frac{P(x)}{Q(x)}$ を考える.分子の次数が分母以上のときは,整式の割り算をして
\[
 P(x)=Q(x)S(x)+R(x)\qquad(R(x)\text{ の次数}<Q(x)\text{ の次数})
\]
と書く.両辺を $Q(x)$ で割ると
\[
 \frac{P(x)}{Q(x)}=S(x)+\frac{R(x)}{Q(x)}
\]
となり,分数全体が「整式」と「分子の次数が低い分数」の和に分かれる.整式 $S(x)$ はべきの公式で積分できる.特に分母が一次式 $x-\alpha$ のときは,余りは定数 $r$ であり,\footnote{剰余の定理より $r=P(\alpha)$ である.割り算の検算に使える.}
\[
 \int\frac{P(x)}{x-\alpha}\dd x=\int S(x)\dd x+r\log|x-\alpha|+C
\]
となる.

\begin{bookpoint}
\textbf{POINT}\par
整式の積は展開し,分子の次数が分母以上の分数は $P=QS+R$ と割って,整式 $S$ と $\frac RQ$ の和に分ける.
\end{bookpoint}

\begin{example}{積を展開する理由}
$\int x(2x-1)^2\dd x$ を求めよ.
\end{example}

\paragraphbf{方針}
2乗を展開して $x$ を掛け,\,$x$ のべきの和にする.

\paragraphbf{解答}
$x^4-\frac43x^3+\frac12x^2+C$

\paragraphbf{解法}
2乗を展開してから $x$ を掛けると
\begin{align*}
 x(2x-1)^2
 &=x(4x^2-4x+1)\\
 &=4x^3-4x^2+x
\end{align*}
となる.したがって
\[
 \int x(2x-1)^2\dd x
 =x^4-\frac43x^3+\frac12x^2+C
\]
である.整式なので,実数全体で成り立つ.右辺を微分すれば $4x^3-4x^2+x$ に戻り,これをくくり直すと元の式になる.

\begin{example}{分子の次数を下げる}
$x\ne-1$ において,\,$\int\frac{x^3+2x^2+3}{x+1}\dd x$ を求めよ.
\end{example}

\paragraphbf{方針}
分子を分母 $x+1$ で割り,整式と $\frac{r}{x+1}$ の和に分ける.

\paragraphbf{解答}
$\frac{x^3}{3}+\frac{x^2}{2}-x+4\log|x+1|+C$

\paragraphbf{解法}
分子を分母で割ると
\[
 x^3+2x^2+3=(x+1)(x^2+x-1)+4
\]
となる.余りは,分子に $x=-1$ を代入した $-1+2+3=4$ と一致する.両辺を $x+1$ で割って
\[
 \frac{x^3+2x^2+3}{x+1}
 =x^2+x-1+\frac4{x+1}
\]
とすれば,すべて既知の形である.よって
\[
 \int\frac{x^3+2x^2+3}{x+1}\dd x
 =\frac{x^3}{3}+\frac{x^2}{2}-x+4\log|x+1|+C
\]
を得る.もとの関数は $x=-1$ で定義されないので,この答えは $x<-1$ と $x>-1$ の各区間で成り立つ.

\paragraphbf{展開が向かない場合}
展開は,数項で済むならば見通しのよい方法である.しかし,高いべきでは展開する項数が多くなるため,別の方法が有利になる.これは本章の終わりで改めて取り上げる.また,割り算を済ませても,余り $\frac{R(x)}{Q(x)}$ の分母が2次以上のことがある.その分数を分けるのが次節の部分分数である.
%
\section{部分分数}\label{節:部分分数}

割り算を終えても,分母が2つの一次式の積になっていることがある.一般的な形は
\[
 \int\frac{px+q}{(x-\alpha)(x-\beta)}\dd x\qquad(\alpha\ne\beta)
\]
で与えられる.分数の分母を勝手に分配することはできない.たとえば $\frac1{x(x+1)}$ は $\frac1x+\frac1{x+1}$ ではない.しかし,一次式の逆数ならば対数で積分できる.そこで,この分数を $\frac A{x-\alpha}$ と $\frac B{x-\beta}$ の和に書き換えて,対数に帰着させることを目標とする.この書き換えを\textbf{部分分数分解}という.

ここで1つ断りをいれておこう.本節の分数は,分子の次数が分母より低いもの,すなわち割り算を済ませたものとする.また,分母を $0$ にする $x=\alpha,\beta$ では定義されないので,答えはそれらの点をまたがない各区間で読む.

右辺を通分すると
\[
 \frac A{x-\alpha}+\frac B{x-\beta}
 =\frac{A(x-\beta)+B(x-\alpha)}{(x-\alpha)(x-\beta)}
\]
となる.したがって,分子どうしの等式
\[
 px+q=A(x-\beta)+B(x-\alpha)
\]
が $x$ についての恒等式となるように $A,B$ を選べばよい.\,$x=\alpha$ を代入すると $B$ の項が消えて
\[
 p\alpha+q=A(\alpha-\beta)
\]
となり,
\[
 A=\frac{p\alpha+q}{\alpha-\beta}
\]
を得る.同様に $x=\beta$ を代入して
\[
 B=\frac{p\beta+q}{\beta-\alpha}
\]
を得る.これは,元の分数に分母を $0$ にする値を代入したのではない.分母を払った整式の恒等式に代入したのである.逆にこの $A,B$ を選べば,両辺の差は1次以下の整式であり,\,$x=\alpha,\beta$ の2点で $0$ になるから,恒等的に $0$ である.よって
\[
 \int\frac{px+q}{(x-\alpha)(x-\beta)}\dd x
 =A\log|x-\alpha|+B\log|x-\beta|+C
\]
となる.

分母に同じ一次因子が重なって $(x-\alpha)^2$ となっている場合は,上の式で $\beta=\alpha$ とすることはできない.\,$A,B$ の分母が $0$ になるからである.この場合は,分子に分母と同じ $x-\alpha$ を作ってやればよい.
\[
 px+q=p(x-\alpha)+(p\alpha+q)
\]
と書けば,
\[
 \frac{px+q}{(x-\alpha)^2}
 =\frac p{x-\alpha}+\frac{p\alpha+q}{(x-\alpha)^2}
\]
となる.第1項は対数で,第2項は $(x-\alpha)^{-2}$ なのでべきの公式で積分でき,
\[
 \int\frac{px+q}{(x-\alpha)^2}\dd x
 =p\log|x-\alpha|-\frac{p\alpha+q}{x-\alpha}+C
\]
を得る.

\begin{bookpoint}
\textbf{POINT}\par
分母の因数ごとに $\frac A{x-\alpha}$ の形へ分け,係数を恒等式で決めて対数で積分する.重なる因子 $(x-\alpha)^2$ は,分子を $x-\alpha$ と定数の和に書いて分ける.
\end{bookpoint}

\begin{example}{部分分数}
$\int\frac{5x+1}{(x-1)(x+2)}\dd x$ を求めよ.
\end{example}

\paragraphbf{方針}
$\frac A{x-1}+\frac B{x+2}$ と置き,分母を払った恒等式に $x=1,-2$ を代入して係数を決める.

\paragraphbf{解答}
$2\log|x-1|+3\log|x+2|+C$

\paragraphbf{解法}
\[
 \frac{5x+1}{(x-1)(x+2)}
 =\frac A{x-1}+\frac B{x+2}
\]
と置き,分母を払うと
\[
 5x+1=A(x+2)+B(x-1)
\]
となる.\,$x=1$ を代入すると $6=3A$,\,$x=-2$ を代入すると $-9=-3B$ なので,\,$A=2$,\,$B=3$ である.したがって
\[
 \int\frac{5x+1}{(x-1)(x+2)}\dd x
 =2\log|x-1|+3\log|x+2|+C
\]
を得る.元の関数は $x=-2,1$ で定義されないので,この答えは $x<-2$,\,$-2<x<1$,\,$1<x$ の各区間で成り立つ.

\begin{example}{重複する因子}
$\int\frac{2x+3}{(x+1)^2}\dd x$ を求めよ.
\end{example}

\paragraphbf{方針}
分子を $2(x+1)+1$ と書き,\,$\frac2{x+1}$ と $\frac1{(x+1)^2}$ に分ける.

\paragraphbf{解答}
$2\log|x+1|-\frac1{x+1}+C$

\paragraphbf{解法}
分母にある $x+1$ を分子にも作りたいので,\,$2x+3=2(x+1)+1$ と書く.すると
\[
 \frac{2x+3}{(x+1)^2}
 =\frac2{x+1}+\frac1{(x+1)^2}
\]
となる.第1項には対数,第2項にはべきの公式を使い,
\[
 \int\frac{2x+3}{(x+1)^2}\dd x
 =2\log|x+1|-\frac1{x+1}+C
\]
を得る.\,$x=-1$ を除き,\,$x<-1$ と $x>-1$ の各区間で成り立つ.

\paragraphbf{係数比較による決め方}
係数は,恒等式の係数を比べて決めることもできる.本節の1つ目の例題では,右辺を整理すると
\[
 5x+1=(A+B)x+(2A-B)
\]
となるから,
\[
 A+B=5,\qquad 2A-B=1
\]
を解いて,同じ $A=2$,\,$B=3$ を得る.係数比較では連立方程式を解くことになるが,\,$x=\alpha$ を代入する方法では,1回の代入で係数が1つずつ決まる.分母の因数が3つ以上に増えても,この代入は同じように使える.

一般に有理関数,すなわち整式どうしの商を積分するときは,まず割り算で分子の次数を下げ,次に分母を因数分解して分けることを考える.ただし,実数の範囲で一次式に分かれない二次式もある.その場合に役立つ形を次節で調べる.
%
\section{平方完成}\label{節:平方完成と積分}

前節では,分母が一次式の積に分かれていた.では,分母や根号の中に二次式 $x^2+bx+c$ があり,因数分解がすぐには見えない場合はどうすればよいだろうか.ここで1つの知識を思い出そう.平方完成である.
\[
 x^2+bx+c=\left(x+\frac b2\right)^2+c-\frac{b^2}4
\]
$m=-\frac b2$,\,$k=c-\frac{b^2}4$ と置けば,この二次式は
\[
 (x-m)^2+k
\]
と書ける.\,$x$ そのものよりも,\,$x-m$ を1つのまとまりとして見ると,二次式の中にある構造が明確になる.本節では,このまとまりを手がかりに,前節の部分分数か,合成関数の微分の逆に帰着させる.

進み方は $k$ の符号で分かれる.\,$k<0$ のときは,\,$k=-h^2$ $(h>0)$ と書けば平方の差になり,
\[
 (x-m)^2-h^2=(x-m-h)(x-m+h)
\]
と因数分解できる.あとは前節の部分分数に帰着する.\,$k=0$ のときは $(x-m)^2$ であり,一次式のべきとして積分できる.

$k>0$ のときは,\,$(x-m)^2+k\ge k>0$ なので,実数の一次式の積には分解できない.このときは,まとまり $(x-m)^2+k$ を微分すると $2(x-m)$ になることを使う.積分される式の中に $x-m$ が掛かっていれば,合成関数の微分を逆に読んで原始関数の候補を作れる.

\begin{bookpoint}
\textbf{POINT}\par
$x^2+bx+c=(x-m)^2+k$ と平方完成し,\,$k<0$ なら因数分解して部分分数へ,\,$k>0$ なら $x-m$ が掛かっているかを見て合成関数の微分を逆に読む.
\end{bookpoint}

\begin{example}{根号の微分と積分}
$\int\frac{x-2}{\sqrt{x^2-4x+5}}\dd x$ を求めよ.
\end{example}

\paragraphbf{方針}
根号の中を $(x-2)^2+1$ と平方完成し,\,$\sqrt{(x-2)^2+1}$ を微分した式と見比べる.

\paragraphbf{解答}
$\sqrt{x^2-4x+5}+C$

\paragraphbf{解法}
平方完成すると,積分される式は
\[
 \frac{x-2}{\sqrt{(x-2)^2+1}}
\]
となる.分子に $x-2$ があり,分母に平方根があるので,\,$\sqrt{(x-2)^2+1}$ を微分してみよう.合成関数の微分から
\begin{align*}
 \frac{\dd}{\dd x}\sqrt{(x-2)^2+1}
 &=\frac1{2\sqrt{(x-2)^2+1}}\cdot2(x-2)\\
 &=\frac{x-2}{\sqrt{(x-2)^2+1}}
\end{align*}
となる.これがちょうど元の式なので,
\[
 \int\frac{x-2}{\sqrt{x^2-4x+5}}\dd x
 =\sqrt{x^2-4x+5}+C
\]
を得る.根号の中は $(x-2)^2+1\ge1$ なので,この答えは実数全体で使える.

\paragraphbf{まとまりの微分がない場合}
同じ根号でも,分子に $x-2$ がない $\int\frac1{\sqrt{(x-2)^2+1}}\dd x$ では,いまの微分の結果と一致しない.その違いを無視して同じ答えにはできない.また,\,$k>0$ で分母が $(x-m)^2+k$ そのものである $\int\frac1{(x-m)^2+k}\dd x$ も,この節の変形だけでは終わらない.前者は\sectionref{節:平方根と部分積分}で,後者は\sectionref{節:三角置換の初歩}の三角置換で扱う.平方完成はいつでも積分を終わらせる操作ではないが,その後に因数分解を使えるのか,置換を考えるべきなのかを判断する材料になるのだ.
%
\section{三角関数の積分}\label{節:三角式の整理}

$\sin x$ と $\cos x$ の原始関数は分かっていても,\,$\sin^2x$ や $\sin(3x)\cos x$ のような2乗や積は,基本公式にそのまま当てはめられない.そこで,三角関数の恒等式を使って $\sin(kx)$ や $\cos(kx)$ の和に直し,\sectionref{節:一次式の積分}の一次式の積分に帰着させよう.

2乗には二倍角の公式
\[
 \cos(2x)=1-2\sin^2x=2\cos^2x-1
\]
を使う.これを $\sin^2x$,\,$\cos^2x$ について解き直して
\[
 \sin^2x=\frac{1-\cos(2x)}2,\qquad
 \cos^2x=\frac{1+\cos(2x)}2
\]
とする.左辺は2乗だが,右辺は定数と余弦の和になった.

積には加法定理を使う.\,$\sin(\alpha+\beta)$ と $\sin(\alpha-\beta)$ の式を並べて足すと,
\[
 \begin{array}{rrl}
  &\sin(\alpha+\beta)={}&\sin\alpha\cos\beta+\cos\alpha\sin\beta\\
  +\,)&\sin(\alpha-\beta)={}&\sin\alpha\cos\beta-\cos\alpha\sin\beta\\ \hline
  &\sin(\alpha+\beta)+\sin(\alpha-\beta)={}&2\sin\alpha\cos\beta
 \end{array}
\]
となり,右辺の第2項が打ち消し合う.よって
\[
 \sin\alpha\cos\beta=\frac12\{\sin(\alpha+\beta)+\sin(\alpha-\beta)\}
\]
を得る.同様に,\,$\cos(\alpha-\beta)$ と $\cos(\alpha+\beta)$ の式を足し引きすると
\[
 \cos\alpha\cos\beta=\frac12\{\cos(\alpha-\beta)+\cos(\alpha+\beta)\},
\]
\[
 \sin\alpha\sin\beta=\frac12\{\cos(\alpha-\beta)-\cos(\alpha+\beta)\}
\]
となる.\,$\alpha=mx$,\,$\beta=nx$ とすれば,積が $\sin$ や $\cos$ の一次式の和に変わる.あとは各項を,角の $x$ の係数で割って積分すればよい.

\begin{bookpoint}
\textbf{POINT}\par
2乗は二倍角の公式で,積は加法定理を2つ足し引きして,\,$\sin(kx)$,\,$\cos(kx)$ の和に直し,\,$k$ で割って積分する.
\end{bookpoint}

\begin{example}{二乗と積を和へ直す}
次の積分を求めよ.
\[
 \text{(1)}\quad\int_0^{\frac\pi2}\sin^2x\dd x,
 \qquad
 \text{(2)}\quad\int\sin(3x)\cos x\dd x.
\]
\end{example}

\paragraphbf{方針}
\begin{enumerate}[label=(\arabic*)]
 \item 二倍角の公式で $\sin^2x=\frac{1-\cos(2x)}2$ と直す.
 \item 積を和に直す式で $\alpha=3x$,\,$\beta=x$ とする.
\end{enumerate}

\paragraphbf{解答}
\begin{enumerate}[label=(\arabic*)]
 \item $\frac\pi4$
 \item $-\frac18\cos(4x)-\frac14\cos(2x)+C$
\end{enumerate}

\paragraphbf{解法}
\begin{enumerate}[label=(\arabic*)]
 \item 二倍角の公式から
 \begin{align*}
  \int_0^{\frac\pi2}\sin^2x\dd x
  &=\frac12\int_0^{\frac\pi2}\{1-\cos(2x)\}\dd x\\
  &=\left[\frac x2-\frac14\sin(2x)\right]_0^{\frac\pi2}\\
  &=\frac\pi4
 \end{align*}
 となる.\,$\cos(2x)$ の積分で $2$ で割るので,全体の係数が $\frac14$ になる.定積分なので,任意定数は上下端で差を取ると消える.
 \item $\sin\alpha\cos\beta=\frac12\{\sin(\alpha+\beta)+\sin(\alpha-\beta)\}$ で $\alpha=3x$,\,$\beta=x$ とすると
 \[
  \sin(3x)\cos x=\frac12\{\sin(4x)+\sin(2x)\}
 \]
 となる.よって
 \[
  \int\sin(3x)\cos x\dd x
  =-\frac18\cos(4x)-\frac14\cos(2x)+C
 \]
 を得る.実数全体で成り立つ.右辺を微分すると $\frac12\sin(4x)+\frac12\sin(2x)$ となり,上の恒等式で積に戻る.
\end{enumerate}

\paragraphbf{公式を忘れたとき}
積を和へ直す3つの式を覚えていなくても,本節の初めのように加法定理を2つ書いて足し引きすれば,必要な式をその場で作れる.正弦と余弦の積なら正弦の加法定理を,正弦どうし,余弦どうしの積なら余弦の加法定理を並べればよい.

\paragraphbf{式の整理の限界}
本章では,積分記号の中の式を同じ値の別の式に書き換えてきた.どの変形を使うかは,元の式に現れた手がかりで決まる.整式の小さな積なら展開し,分子の次数が高ければ割り算をし,一次因子の積が分母にあれば部分分数を考える.二次式は平方完成でまとまりを作り,三角関数の2乗や積は恒等式で和へ直す.

ただし,式を長く展開すれば必ずよいわけではない.たとえば $x(x^2+1)^{20}$ は展開できるが,\,$x^2+1$ を1つの変数 $t$ と置けば計算を短くできる.そう置いてよい根拠を与えるのが,次章で学ぶ置換積分である.
%
\needspace{18\baselineskip}
\section*{章末問題}
\begin{enumerate}[label=\textbf{\thechapter.\arabic*},leftmargin=*]
 \item\label{問:基本べき} $x>0$ において,\,$\int\frac{3x^2-2x+4}{\sqrt{x}}\dd x$ を求め,微分して確かめよ.
 \item\label{問:基本割算} $\int\frac{x^3+1}{x+1}\dd x$ を求めよ.式を整理した後も,もとの関数で除かれる点を述べよ.
 \item\label{問:基本部分分数} $\int\frac{5x-1}{(x-1)(x+1)}\dd x$ を求めよ.部分分数の係数を決める等式も書け.
 \item\label{問:基本重因子} $\int\frac{3x+4}{(x+1)^2}\dd x$ を求めよ.
 \item\label{問:基本平方完成} 分母を平方完成してから,\,$\int\frac1{x^2+6x+8}\dd x$ を求めよ.
 \item\label{問:基本三角二乗} $\int_0^{\frac\pi2}\cos^2(2x)\dd x$ を求めよ.
 \item\label{問:基本三角積} 加法定理から必要な積を和へ直す式を導き,\,$\int\sin(2x)\sin(5x)\dd x$ を求めよ.
\end{enumerate}

\smallskip\noindent 解答は\pageref{解答:第5章}ページ.

%
